\documentclass[../../../hsm_tok_q.tex]{subfiles}

\begin{document}
    \setcounter{figure}{0}
    \textsf{図\ref{fig/hsm_tok_2025_2nd_04_01_q}}で，%
    四角形ABCDは，平行四辺形である．

    点Pは，辺BC上にある点で，頂点B，頂点Cのいずれにも一致しない．

    頂点Aと点Pを結ぶ．
    
    次の各問に答えよ．
    \\
    \begin{minipage}{\linewidth}
        \vspace{.5\baselineskip}
        {\centering
            \setlength{\abovecaptionskip}{3pt}
            \includegraphics%
                {\subfix{fig_hsm_tok_2025_2nd_04/fig_hsm_tok_2025_2nd_04_01_q.pdf}}%
                \captionof{figure}{} \label{fig/hsm_tok_2025_2nd_04_01_q}%
        }%
    \end{minipage}
    \vspace{.5\baselineskip}

    \begin{enumerate}
        \item \textsf{図\ref{fig/hsm_tok_2025_2nd_04_01_q}}において，%
            $\angle\mathrm{ABC}\jequal50^\circ$，$\angle\mathrm{BAP}\jequal a^\circ$ とするとき，%
            $\angle\mathrm{DAP}$の大きさを表す式を，%
            次の\textsf{ア}～\textsf{エ}のうちから選び，記号で答えよ．

            {\renewcommand{\arraystretch}{1.3}%
                \vspace{-.2\baselineskip}
                \begin{tabularx}{\dimexpr\linewidth-3.5mm}{@{}LLLL@{}}
                    {\textsf ア}\; $(50-a)$度 & {\textsf イ}\; $(a+50)$度 \\
                    {\textsf ウ}\; $(130-a)$度 & {\textsf エ}\; $(a+130)$度
                \end{tabularx}
            }
        \newpage
        \item \textsf{図\ref{fig/hsm_tok_2025_2nd_04_02_q}}は，%
            \textsf{図\ref{fig/hsm_tok_2025_2nd_04_01_q}}において，%
            頂点Bと頂点Dを結び，%
            線分APと線分BDとの交点をQとした場合を表している．

            次の\ref{hsm_tok_2025_2nd_04_q:1}，\ref{hsm_tok_2025_2nd_04_q:2}に答えよ．
            %
            \begin{myenub}
                \item $\triangle\mathrm{AQD}\jsim\triangle\mathrm{PQB}$であることを証明せよ．
                    \label{hsm_tok_2025_2nd_04_q:1}

                \item 次の\:\myboxf{　}\:の中の「\textsf{う}」「\textsf{え}」「\textsf{お}」%
                    に当てはまる数字をそれぞれ答えよ．

                    \textsf{図\ref{fig/hsm_tok_2025_2nd_04_02_q}}において，%
                    点Pを通り線分BDと平行な直線を引き，辺CDとの交点をR，頂点Aと点Rを結び，%
                    線分ARと線分BDとの交点をSとし，点Pと点Sを結んだ場合を考える．

                    $\mathrm{BP}\jcolon\mathrm{PC}\jequal3\jcolon2$ のとき，%
                    $\triangle\mathrm{PSQ}$の面積は，四角形ABCDの面積の %
                    $\myfracc{\myboxf{う}}{\myboxf{えお}}$ 倍である．
                    \label{hsm_tok_2025_2nd_04_q:2}
            \end{myenub}
            %
            \begin{minipage}{\linewidth}
                \vspace{.5\baselineskip}
                {\centering
                    \setlength{\abovecaptionskip}{3pt}
                    \includegraphics%
                        {\subfix{fig_hsm_tok_2025_2nd_04/fig_hsm_tok_2025_2nd_04_02_q.pdf}}%
                        \captionof{figure}{} \label{fig/hsm_tok_2025_2nd_04_02_q}%
                }%
            \end{minipage}            
        \end{enumerate}
\end{document}

